\subsection{Hypothesis and Test Procedures}
\hrulefill

\subsubsection*{Hypothesis Testing}
\paragraph*{Definition}
A statistical hypothesis is a claim about a population parameter. The null hypothesis, denoted by $H_0$, is the claim that is initially
assumed to be true. The alternative hypothesis, denoted by $H_a$, is an assertion that is contradictory to $H_0$. Null hypothesis is 
the status quo, aka the default assumption is absence of compelling evidence. The alternative hypothesis is what we want to find evidence for.

At the conclusion of a hypothesis test, we either reject $H_0$ in favor of $H_a$ or fail to reject $H_0$. We never accept $H_0$ as true, 
we only fail to reject it. An analogue to this is a court trial, where the defendant is assumed innocent until proven guilty. 
If there is enough evidence to convict, the defendant is found guilty (reject $H_0$). If there is not enough evidence, the defendant is
found not guilty (fail to reject $H_0$).

\paragraph*{Errors}
There are two possible errors that can be made in hypothesis testing:
\begin{itemize}
    \item Type I Error: Rejecting the null hypothesis when it is actually true. (Finding a guilty defendant innocent)
    \item Type II Error: Failing to reject the null hypothesis when it is actually false. (Finding an innocent defendant guilty)
\end{itemize}

\begin{itemize}
    \item We want to keep these two types of error rates as small as possible
    \item However, it is impossible to minimize both simultaneously
    \item The choice of rejection region determines the probabilities of both a type I and type II error
\end{itemize}

\paragraph*{Significance Level}
The probability of a type I error is called the significance level of the test, denoted by $\alpha$. The probability of a type II error 
is usually denoted by $\beta$. We agree on a maximum acceptable value for $\alpha$ before conducting the test, commonly 0.05, 0.01, or 0.005.
The quantity $1 - \beta$ is called the power of the test, which is the probability of correctly rejecting a false null hypothesis. Good tests
have high power.

\paragraph*{Components of a Statistical Test}
\begin{itemize}
    \item A \textbf{test statistic} (based on a random sample from the population under question) that summarized the evidence against the null hypothesis
    \item A \textbf{p-value} is a probability that measures how likely the value of the test statistic is, assuming that $H_0$ is true
    \item A \textbf{rejection region} tells us for which values of the test statistic we reject $H_0$
    \item A \textbf{Conclusion} is a decision of the form either Reject $H_0$ or Fail to Reject $H_0$
    \item A \textbf{significance level} $\alpha$ which represents the maximum allowable risk of rejection $H_0$ in favor of $H_a$ when $H_0$ is true.
\end{itemize}

\paragraph*{p-values}
One way to report the results of a hypothesis test is to determine if the test statistic value falls into the rejection region and subsequently, whether or not the 
null hypothesis is rejected at a specified level of significance. This yes/no decision does not convey any information about how soundly the 
null hypothesis was rejected. Where in the rejection region did the observed test statistic fall? Note that the p-value is calculated assuming
that the null hypothesis is true.